\chapter{Conclusioni e Sviluppi Futuri}
\label{chap:conclusioni}

Il lavoro svolto in questa tesi ha affrontato il problema della raccolta 
automatizzata di pomodori in serra, concentrandosi sulla pianificazione 
della sequenza di prelievo all'interno di grappoli densi e complessi.

Nel quadro di un approccio di visione attiva (\textit{Active Vision}), l'analisi accurata 
della singola inquadratura rappresenta il primo passo indispensabile: prima ancora di 
muovere la telecamera attorno alla pianta per cercare frutti nascosti, il sistema deve 
saper gestire le occlusioni visibili e pianificare una sequenza sicura per evitare collisioni.

Di seguito vengono riassunti i principali risultati raggiunti con la 
modellazione ellittica e la ripianificazione dinamica, delineando poi 
le prospettive per i futuri sviluppi e l'integrazione robotica.


\section{Sintesi dei Risultati e Contributi della Tesi}
\label{sec:sintesi_risultati}

La raccolta robotizzata in serra richiede di individuare con precisione 
l'ordine di presa dei frutti, soprattutto all'interno di grappoli densi dove 
le occlusioni reciproche rendono rischioso l'avvicinamento del robot. I metodi 
tradizionali, basati su riquadri rettangolari rigidi (\textit{bounding box}) o stime statiche, 
mostrano forti limiti quando i pomodori sono a stretto contatto o presentano forme allungate.

Per superare queste criticità, è stata sviluppata e validata una pipeline 
decisionale basata su tre passaggi principali:
\begin{enumerate}
    \item \textbf{Regolarizzazione geometrica con Convex Hull}: l'applicazione 
    dell'involucro convesso sulle maschere di YOLO11-x ha permesso di ripulire i 
    bordi a gradino, rendendo il calcolo del perimetro e della circolarità stabile 
    e realistico;
    \item \textbf{Modellazione ellittica orientata}: l'approssimazione dei pomodori 
    mediante ellissi orientate ha eliminato le false occlusioni d'angolo tipiche 
    dei rettangoli rigidi (in cui l'area vuota negli angoli occupa circa il $21.5\%$ 
    della scatola), adattandosi all'inclinazione e alla forma reale del frutto;
    \item \textbf{Task planning dinamico \textit{Pick \& Update}}: la simulazione 
    del prelievo a ogni passo con l'assegnazione del bonus di sblocco 
    ($\beta_{\text{unlock}} = 1.20$) ha consentito di ripianificare la sequenza, 
    premiando i frutti retrostanti appena liberati e trattando i pomodori acerbi 
    come ostacoli statici da non urtare.
\end{enumerate}

La sperimentazione condotta sulle immagini reali di serra ha confermato la validità del 
metodo proposto. 

Analizzando l'intero set di $43$ scene, nel $62.8\%$ dei casi ($27$ scene) la sequenza 
prodotta coincide con quella della baseline, confermando che nei grappoli semplici il 
modello non introduce distorsioni. Nelle scene dense ($12$ casi, $27.9\%$), l'algoritmo 
corregge errori di precedenza dovuti a contatti diagonali, mentre in $4$ scene ($9.3\%$) 
rileva correttamente l'assenza di frutti maturi accessibili, arrestando la raccolta 
in sicurezza (condizione di \textit{safe-fail}).

Sul benchmark di $21$ scene ad alta densità con riferimento umano, sia la baseline sia 
il modello proposto registrano una concordanza dell'$81.0\%$ sul primo frutto da raccogliere 
(Rank~1, $17$ casi su $21$). Tuttavia, un'analisi puntuale delle scene rivela come i due 
approcci giungano a tale risultato per vie profondamente diverse. Nelle situazioni in cui 
la baseline fallisce, l'errore è causato dalla geometria rigida dei riquadri rettangolari, 
i cui angoli vuoti inducono false occlusioni che spingono il sistema a selezionare frutti 
parzialmente coperti. Al contrario, le divergenze del planner ellittico non comportano 
rischi di collisione, ma si verificano quando più pomodori risultano ugualmente accessibili: 
in tali contesti, l'algoritmo privilegia deterministicamente il frutto con il punteggio 
geometrico più elevato (ad esempio per una maggiore centralità ottica), a fronte di una 
scelta soggettiva dell'operatore.

La reale efficacia della modellazione curvilinea e della pianificazione dinamica emerge 
sui prelievi successivi. Grazie all'aggiornamento iterativo dello stato e all'introduzione 
del bonus di sblocco ($\beta_{\text{unlock}} = 1.20$), la concordanza con l'operatore umano 
sale al $46.2\%$ sul secondo frutto (rispetto al $38.5\%$ della baseline) e raggiunge il 
$40.0\%$ sul terzo, raddoppiando l'accuratezza del modello convenzionale ($20.0\%$).

Infine, dal punto di vista computazionale, la scelta di selezionare a monte un pool 
ristretto di $K = 8$ candidati limita i confronti tra ellissi a un massimo di $56$ 
verifiche geometriche. Questo rende l'algoritmo leggero e veloce, adatto a essere 
eseguito in tempo reale a bordo di un sistema robotico.


\section{Sviluppi Futuri e Integrazione Robotica}
\label{sec:sviluppi_futuri}

I risultati ottenuti in questo lavoro confermano l'efficacia dell'approccio geometrico 
e dinamico per la pianificazione della raccolta su singola inquadratura. Trattandosi di 
uno studio incentrato sulla componente algoritmica e decisionale, il naturale proseguimento 
di questa ricerca riguarda il trasferimento della pipeline a bordo di un manipolatore 
robotico operativo e la risoluzione dei limiti pratici riscontrati durante la sperimentazione.

Il primo passo verso l'applicazione fisica consiste nel passaggio dalle coordinate del piano 
immagine allo spazio operativo tridimensionale. Utilizzando una telecamera di profondità 
RGB-D, la posizione bidimensionale $(u, v)$ del centroide calcolato dalla nostra pipeline 
può essere combinata con il canale di profondità per ricavare le coordinate cartesiane $(X, Y, Z)$ 
nel sistema di riferimento ottico della telecamera, mediante la matrice dei parametri intrinseci 
del sensore \cite{gongal2015sensors}. In questa fase, per evitare che rumore ottico o riflessi 
alterino la misura, la profondità $Z$ del pomodoro potrà essere calcolata come mediana dei pixel 
validi all'interno della maschera di segmentazione. Inoltre, l'angolo di inclinazione $\theta$ 
restituito dal fitting dell'ellisse fornirà l'orientamento necessario per allineare la pinza di presa 
lungo l'asse principale del frutto, agevolando il distacco meccanico dal peduncolo. Attraverso il 
middleware ROS~2 e il pacchetto di trasformazioni spaziali \texttt{tf2}, le coordinate verranno 
quindi convertite rispetto alla base del manipolatore (come il braccio industriale ABB presente 
nel setup sperimentale) e inviate a MoveIt~2 \cite{chitta2012moveit}. Il framework di motion planning 
potrà così generare traiettorie prive di collisioni nello spazio dei giunti, guidando l'organo di presa 
verso i frutti secondo l'ordine di priorità calcolato dal nostro planner.

Un secondo fronte di evoluzione riguarda il superamento dei limiti imposti dalla visione 
da un singolo punto di vista. Come osservato analizzando le immagini di test in 
\textit{immaginiPerAlgoritmo}, l'algoritmo gestisce con successo le occlusioni parziali tra 
pomodori adiacenti, ma non può individuare frutti maturi completamente nascosti dietro il fusto 
o sotto uno strato fitto di vegetazione. La soluzione naturale consiste nell'integrare la pipeline 
sviluppata all'interno di una strategia di visione attiva (\textit{Active Vision}) \cite{silwal2017design}. 
Montando il sensore visivo direttamente sul polso del manipolatore (\textit{eye-in-hand}), il robot 
potrà compiere movimenti esplorativi attorno alla pianta per acquisire viste complementari da angolazioni 
diverse (\textit{Next-Best-View}). In questo scenario, la nostra pipeline agirà come decisore locale: a ogni 
nuova inquadratura, il modello mapperà le occlusioni visibili e aggiornerà la lista dei frutti da raccogliere, 
integrando progressivamente la conoscenza del grappolo man mano che i movimenti del braccio svelano i 
pomodori precedentemente occlusi.

Infine, un'ulteriore linea di sviluppo riguarda la robustezza del filtro di maturità. L'analisi 
delle condizioni critiche ha evidenziato come l'impiego di soglie rigide nello spazio colore HSV 
($S > 50, V > 50$) mostri fragilità in presenza di forti zone d'ombra, dove la scarsa luminosità 
rischia di far escludere dal conteggio molti pixel rossi. Per superare questa limitazione senza 
compromettere la semplicità computazionale del sistema, si prospettano due alternative: la transizione 
allo spazio colore CIELAB ($L^*a^*b^*$), in cui il canale cromatico $a^*$ (che misura la transizione dal 
verde al rosso) è disaccoppiato dall'intensità luminosa $L^*$ risultando pressoché immune alle ombre; 
oppure l'addestramento di una rete convoluzionale leggera (come MobileNetV3) operante direttamente 
sui ritagli dei singoli pomodori individuati da YOLO \cite{xiao2023fruit}. Quest'ultima soluzione 
permetterebbe di stimare il grado di maturazione basandosi anche sulla tessitura superficiale del frutto, 
riconoscendo con maggiore affidabilità anche le fasi intermedie di viraggio del colore.

